\documentclass{article}

\usepackage{amsmath}
\usepackage{amssymb}
\usepackage{geometry}
\geometry{margin=2.5cm}
\title{TOX Flyer coding guidelines}
\author{Asis Hallab, ``el Bobito''}

\begin{document}

\maketitle

\section{Rendering Abstraction in the TOX Flyer}

This section introduces a light-weight rendering abstraction strategy for the TOX Flyer, based on the current Babylon.js implementation. The aim is not to introduce unnecessary complexity, but to gently decouple rendering logic from application logic in a way that respects the KISS principle and anticipates future portability needs (e.g., migration to WebGPU).

\subsection{Current Implementation Pattern}

In the present implementation, Babylon-specific drawing commands are used directly inside the main application logic. For instance:

\begin{verbatim}
const sphere = BABYLON.MeshBuilder.CreateSphere("s", {diameter: 1}, scene);
sphere.position = new BABYLON.Vector3(x, y, z);
sphere.material = new BABYLON.StandardMaterial("mat", scene);
sphere.material.diffuseColor = new BABYLON.Color3(1, 0, 0);
\end{verbatim}

This pattern works well for fast development but makes it difficult to:
\begin{itemize}
  \item Replace Babylon.js with another engine in the future (e.g., raw WebGPU).
  \item Test scene logic independently of Babylon’s rendering behavior.
  \item Provide reusable components across the Flyer or future viewers.
\end{itemize}

\subsection{Proposed Abstraction Pattern}

The following pattern introduces a thin, declarative rendering interface that encapsulates Babylon-specific logic while remaining simple and readable.

\subsubsection*{Example: Minimal Rendering Interface}

\texttt{render\_interface.js}:
\begin{verbatim}
export function addSphere(scene, position, color, label) {
  const sphere = BABYLON.MeshBuilder.CreateSphere(label, {diameter: 1}, scene);
  sphere.position = new BABYLON.Vector3(...position);
  sphere.material = new BABYLON.StandardMaterial("mat", scene);
  sphere.material.diffuseColor = BABYLON.Color3.FromHexString(color);
  return sphere;
}
\end{verbatim}

\subsubsection*{Usage in Application Code}

\begin{verbatim}
import { addSphere } from "./render_interface.js";

const s = addSphere(scene, [x, y, z], "#ff0000", "module_17");
\end{verbatim}

\subsubsection*{Benefits}

\begin{itemize}
  \item Babylon logic is isolated in one file.
  \item Scene logic becomes declarative and easier to read.
  \item Future refactoring (e.g., to WebGPU or three.js) only requires changes in the interface layer.
\end{itemize}

\subsection{Scope and Philosophy}

This rendering abstraction is:
\begin{itemize}
  \item \textbf{Not a framework}: No additional logic or state management is introduced.
  \item \textbf{Not a rewrite}: Existing drawing commands are simply wrapped.
  \item \textbf{Not a barrier to speed}: Developers may still use Babylon directly when rapid testing is needed.
\end{itemize}

It is a lightweight buffer between creative experimentation and technical future-proofing — aligned with F42 principles and designed to scale with the increasing complexity of the TOX Flyer.

\subsection{Conclusion}

Adopting this minimal RAL allows the TOX Flyer to remain agile while building a foundation for long-term reuse and interoperability. It ensures the Flyer can later be re-targeted (e.g., to WebGPU) without re-engineering all application logic, while preserving the clarity and joy of its current Babylon-based design.

\end{document}
